\section{Methods}
\label{sec:methods}

Figure~\ref{fig:workflow} summarises the analysis workflow, which was applied separately to each dataset.

\begin{figure}[t]
\centering
\includegraphics[width=\textwidth]{figures/fig2_workflow.pdf}
\caption{Study workflow. For each held-out participant $p$, feature standardisation, model fitting and electrode ranking use only the remaining participants $R_{-p}$. Attributions are computed on participant-disjoint inner folds of $R_{-p}$, and models restricted to the top-ranked $k$ electrodes are retrained on $R_{-p}$ and tested on $p$. The same workflow was applied separately to each dataset.}
\label{fig:workflow}
\end{figure}

\subsection{Terminology}
An \emph{electrode} is a physical recording site, and a \emph{channel} is the signal it provides to the classifier; in the monopolar analyses each electrode yields one channel, so selecting electrodes and selecting channels are the same operation. The \emph{montage} is the spatial arrangement of the electrodes and the \emph{array} is the complete set. A \emph{bipolar pair} is the difference between two electrodes, derived from the monopolar recordings in one exploratory analysis. The subset size $k$ always counts the selection unit of the analysis in question (electrodes, or bipolar pairs).

\subsection{Datasets}

\paragraph{GRABMyo (primary)}
GRABMyo v1.1.0 \citep{pradhan2022multiday,jiang2024grabmyo}, distributed through PhysioNet \citep{goldberger2000physionet}, contains recordings from \GRABMyoParticipants\ able-bodied participants (\GRABMyoMale\ male, \GRABMyoFemale\ female; age \GRABMyoAgeMean\ $\pm$ \GRABMyoAgeSD\ years) on \GRABMyoSessions\ different days. Signals were recorded with an EMGUSB2+ amplifier (OT Bioelettronica) at 2048~Hz. In each session, participants performed 16 hand and wrist gestures and rest, with 7 trials of 5~s each. We analysed the 16 forearm electrodes. These form two circumferential rings of 8 electrodes (F1--F8 and F9--F16) 2~cm apart along the forearm, with electrode $j$ of one ring aligned with electrode $j+8$ of the other. Mean forearm circumference at the ring position was \GRABMyoCircMean\ $\pm$ \GRABMyoCircSD~cm, giving about \GRABMyoArcSpacing~cm between neighbouring electrodes in a ring. The recordings are monopolar, so the selection unit in the primary analysis is a single electrode. All \GRABMyoFilesVerified\ distributed files matched the dataset's published SHA-256 checksums, and all \GRABMyoRecords\ records were analysed. The primary task is the 16 active gestures (rest excluded), for which chance Macro-F1 is about 0.063.

\paragraph{Replication datasets}
Two further datasets were analysed separately, never pooled with GRABMyo, to examine whether the procedure behaves similarly with different hardware and tasks. Their channel numbers have no anatomical correspondence with those of GRABMyo, so only the behaviour of the procedure is compared across datasets.
NinaPro DB5 \citep{pizzolato2017comparison} provides \DBFiveSubjects\ able-bodied subjects (\DBFiveMale\ male, \DBFiveFemale\ female; age \DBFiveAgeMean\ $\pm$ \DBFiveAgeSD\ years) recorded with two Myo armbands (16 channels, 200~Hz). We used exercise~1 (12 finger movements, 6 repetitions each), segmented with the refined movement and repetition labels supplied with the dataset. The DB5 values are 8-bit integer codes, not physical units; their spectra show no powerline peak and extend to the Nyquist frequency.
putEMG \citep{kaczmarek2019putemg} provides 44 able-bodied participants (8 female, 36 male; aged 19--37 years) recorded on two days at least one week apart. Twenty-four monopolar electrodes form three circumferential bands of 8 (electrodes 1--8, 9--16 and 17--24) on the forearm, sampled at 5120~Hz with 12-bit resolution. Each session contains 20 repetitions of each of 7 active gestures, labelled from video-verified ground truth. The putEMG analysis was not part of the original plan; its design was fixed before any putEMG model was trained.

\paragraph{Ethics statement}
This study is a secondary analysis of previously published, de-identified datasets and involved no new participant recruitment or data collection. The original data collections were conducted in accordance with the Declaration of Helsinki and approved by the Office of Research Ethics of the University of Waterloo (ORE\#~31346; GRABMyo), the Ethics Commission of the Canton of Valais, Switzerland (NinaPro DB5) and the Bioethical Committee of Poznan University of Medical Sciences (no.~398/17; putEMG); informed-consent procedures are described in the dataset publications \citep{pradhan2022multiday,pizzolato2017comparison,kaczmarek2019putemg}. [AUTHOR INPUT REQUIRED: state whether an institutional review board determination or exemption was obtained for the present secondary analysis.]

\paragraph{Sex as a reported variable}
Sex was taken as reported by the source datasets (GRABMyo \GRABMyoMale\ male and \GRABMyoFemale\ female; DB5 \DBFiveMale\ male and \DBFiveFemale\ female; putEMG 36 male and 8 female). Sex-stratified analyses were not prespecified, and the study was not designed to estimate sex differences in recognition performance or electrode ranking; no sex-specific inference is made.

\subsection{Preprocessing and features}
\label{sec:preprocessing}
GRABMyo signals were band-pass filtered (20--450~Hz, 4th-order zero-phase Butterworth) and notch filtered at 50~Hz ($Q = 30$), the mains frequency identified from the spectra (mean band-power ratio 10.35 at 50~Hz and 0.62 at 60~Hz). Because the filters are zero-phase, the processing is non-causal and suitable only for offline evaluation. Windows were 250~ms with 50\% overlap and never crossed trial boundaries. From each channel we extracted ten features: mean absolute value (MAV), root mean square, variance, waveform length, zero crossings (ZC), slope sign changes (SSC), integrated EMG (IEMG), Willison amplitude (WAMP), mean frequency and median frequency. This gives a 160-dimensional vector per window: \GRABMyoWindows\ windows in total and \GRABMyoActiveWindows\ from active gestures.

The amplitude threshold of the ZC, SSC and WAMP features was set by one rule for every representation: 1.15 times the median absolute first difference of the signal, rounded to one significant figure. For GRABMyo it was computed once, from 54 unfiltered recordings of 6 participants, giving $10^{-2}$~mV; it was therefore not estimated within each outer fold. Recomputing it with each held-out participant excluded gave the same value in every fold (supplementary methods~S2), so fold-wise estimation would have produced identical features.
DB5 was band-pass filtered at 10--95~Hz without a notch filter, with 250~ms windows (50 samples) and a threshold of 7 quantisation steps from the same rule. putEMG was converted to millivolts and processed as GRABMyo (20--450~Hz band-pass, 50~Hz notch, 250~ms windows with 50\% overlap), using only windows that lay entirely within one gesture; because raw putEMG differences are quantised at one converter step, the rule was applied to the filtered signal and gave \PutThresh${}\times 10^{-3}$~mV, for 240 features per window. For both datasets the leave-one-out recomputation gave the same threshold in every fold.

\subsection{Classifiers and outer leave-one-subject-out evaluation}
All evaluation is leave-one-subject-out (LOSO), with one fold per participant. In fold $p$, standardisation statistics, model fitting and electrode ranking use only $R_{-p}$, and the test participant contributes only to the score. XGBoost was designated in the analysis plan as the primary classifier for the electrode-selection analyses because its computational cost made the required number of retrained models feasible; the plan was written before the full-cohort classifier comparison had finished, although results from an earlier, superseded run on 15 participants were available. XGBoost used 300 trees, depth 6, learning rate 0.1, row and column subsampling of 0.8, histogram-based tree construction and a multiclass softmax objective, and was trained on all \FoldTrainingWindows\ training windows of each fold. Two further classifiers were compared with all electrodes: a support vector machine with a radial basis function kernel ($C = 1$, kernel width set from the feature variance), trained on a stratified sample of 40,000 windows, and a random forest (200 trees, minimum leaf size 2) trained on 200,000 windows. These sample sizes were set by computational cost and chosen on two folds (supplementary methods~S1). Halving the support vector machine's training sample on one fold lowered Macro-F1 from \SVMCalibFortyK\ to \SVMCalibTwentyK, so its result is likely a lower bound. Hyperparameters were fixed in advance and not tuned. The primary metric is participant-level Macro-F1 over the 16 active gestures; accuracy and balanced accuracy were also recorded.

\subsection{Cross-fitted attribution and the SR-CFA score}
Within each outer fold, $R_{-p}$ was split into $K = 5$ participant-disjoint inner folds. For inner fold $j$, an XGBoost model was trained on the other inner folds, using a random sample of at most 200,000 windows standardised with inner-training statistics only, and exact TreeSHAP attributions \citep{lundberg2020local} were computed for up to 20,000 windows sampled at random from the inner-held-out participants. Because participants contribute similar numbers of windows (median \TestWindowsPerParticipant), each participant carries roughly equal weight. The importance of channel $c$ for window $i$ averaged the absolute attributions over the $G = 16$ classes and the $F = 10$ features of the channel,
\begin{equation}
\Phi_{i,c} = \frac{1}{G F}\sum_{g=1}^{G}\sum_{f=1}^{F} \lvert \phi_{i,g,c,f} \rvert ,
\end{equation}
and was averaged over windows and normalised to sum to one across channels, giving $I_{j,c}$. Across the five inner folds, $A_c = \operatorname{median}_j I_{j,c}$ measures attribution magnitude and $U_c = \operatorname{median}_j \lvert I_{j,c} - A_c \rvert$ its instability across disjoint groups of training participants. Electrodes were ranked by
\begin{equation}
Q_c(\lambda) = A_c - \lambda\, U_c , \qquad \lambda \in \{0, 0.25, 0.5, 1\},
\label{eq:srcfa}
\end{equation}
with ties broken by the lower channel number. The value $\lambda = 0$ gives cross-fitted SHAP magnitude and $\lambda > 0$ the stability-regularised score (SR-CFA). In every fold we verified that no inner-held-out participant contributed to the corresponding inner training set. The plan did not specify a criterion for choosing $\lambda$, so all four values are reported.

\subsection{Electrode-count performance curve and non-inferiority (H1)}
For every outer fold, $\lambda$ and $k \in \{1,\ldots,16\}$, XGBoost was retrained from scratch on the top-$k$ electrodes of that fold's ranking, using a fixed, stratified sample of 200,000 training windows of the fold, and tested on the held-out participant. The $k = 16$ model of the same curve is the reference, so each paired difference $\Delta_{p,k} = \mathrm{F1}_{p,k} - \mathrm{F1}_{p,16}$ compares models trained on identical data. Subset size $k$ is non-inferior at margin $\delta$ when $L_k$, the lower bound of the 95\% participant-cluster bootstrap confidence interval of the mean difference $\bar\Delta_k$, exceeds $-\delta$. The smallest non-inferior subset is $k^* = \min\{k < 16 : L_k > -\delta\}$, recorded as `none' if no $k$ qualifies. As required by the plan, the family of 15 tests was also corrected by the Holm procedure, using one-sided cluster-bootstrap $p$-values for $H_0\colon \bar\Delta_k \le -\delta$ at $\alpha = 0.025$.

The primary margin $\delta = 0.02$ is an absolute difference in Macro-F1 on a 0--1 scale. It was fixed before the curve was computed as the largest loss we would accept in exchange for fewer electrodes. Relative to the full-array reference it corresponds to about \MarginPctOfRef\% of performance and about \MarginPctAboveChance\% of the distance above chance (1/16). It is an engineering tolerance and has not been validated against clinical outcomes or user performance. Because $k^*$ depends strongly on $\delta$, the prespecified sensitivity margins $\delta = 0.01$ and $0.03$ are also reported. Non-inferiority within this tolerance does not imply equivalence.

\subsection{Comparison with random subsets and other selection methods (H2)}
At $k \in \{4, 6, 8, 12\}$ we compared six rankings, each computed on the training participants only:
\begin{itemize}
\item SR-CFA ($\lambda = 0.5$);
\item cross-fitted SHAP ($\lambda = 0$);
\item permutation importance: a model was fitted on 20,000 training windows and each of the 160 features was permuted in 8,000 further windows (3 repeats, Macro-F1 scoring); an electrode's score is the mean over its 10 features;
\item mutual information between each feature and the gesture label on the same 20,000 windows, averaged over the electrode's 10 features;
\item geometric selection: farthest-point sampling on the $2 \times 8$ montage (arc distance around a 24~cm circumference, 2~cm between rings, starting from F1, ties to the lower index), which uses no data;
\item random subsets: 100 independent random subsets per participant at each $k$, and 500 at the prespecified $k^*$ (seed fixed per participant).
\end{itemize}
Evaluating hundreds of random subsets per participant required a lighter comparison model (XGBoost with 50 trees, trained on a stratified sample of 150,000 windows), applied identically to every method. Its absolute scores are lower than those of the performance curve and are compared only with each other. The confirmatory H2 family contrasts each attribution ranking (SR-CFA and cross-fitted SHAP) with the participant's mean over its random subsets (Holm over 8 tests). Comparisons with permutation importance, mutual information and geometric selection are exploratory and uncorrected. We also report the percentile of the SR-CFA subset within each participant's random subsets. Below, `SR-CFA' refers to $\lambda = 0.5$ and `cross-fitted SHAP' to $\lambda = 0$.

\subsection{Stability regularisation (H3) and faithfulness (H4)}
H3 compared each $\lambda > 0$ with $\lambda = 0$ on the performance curve at $k \in \{4, 6, 8, 12\}$ (Holm over 12 tests). For H4 we used deletion with retraining \citep{hooker2019benchmark}. For $n \in \{2, 4, \ldots, 14\}$ electrodes retained, the comparison model was retrained after removing the $16 - n$ highest-ranked electrodes (Remove-Top), the $16 - n$ lowest-ranked electrodes (Remove-Bottom), or random electrodes (Remove-Random, the random subsets above, available for $n \in \{4, 6, 8, 12\}$). A faithful ranking gives Remove-Top $<$ Remove-Random $<$ Remove-Bottom. Paired contrasts were Holm-corrected within H4.

\subsection{Statistics}
The participant is the statistical unit ($n = 43$ for GRABMyo, 10 for DB5 and 44 for putEMG), and windows are never treated as independent. Confidence intervals (CI) are 95\% percentile intervals from 10,000 cluster-bootstrap resamples of participants (seed 20260728). Paired comparisons use the Wilcoxon signed-rank test with the matched-pairs rank-biserial correlation $r_{rb}$ as effect size. Holm correction is applied within each confirmatory family; exploratory analyses are uncorrected and labelled as such. A non-significant difference is reported as `no statistically detectable difference', never as equivalence, because no equivalence test was performed. No participant was excluded.

\subsection{Exploratory analyses}
\paragraph{Ranking stability} We compared rankings derived from single participants' attributions (all 903 pairs), from the five inner folds and from the 43 outer folds, using Spearman's $\rho$, Kendall's $\tau$ and the Jaccard index of the top-$k$ sets. A participant's attribution profile is the mean of their normalised attribution over the outer folds in which they were held out of an inner model. Outer folds share about 98\% of their training participants, so their agreement is only an upper bound.

\paragraph{Replication} The LOSO baseline, cross-fitted ranking ($K = 5$) and performance curve were repeated on DB5 and putEMG, each with its full array as the reference ($k = 16$ and $k = 24$), using all training windows and the same bootstrap, margins and $k^*$ rule, with Holm-corrected results also reported. On putEMG, ranked subsets were also compared with ten random subsets per participant at $k \in \{6, 12, 18\}$ (Wilcoxon, Holm over 3). No setting was tuned on either dataset.

\paragraph{Cross-session} XGBoost was trained on session~1 of $R_{-p}$ and tested on each session of $p$, and trained on sessions 1 and 2 and tested on session~3. The second condition also has about twice the training data, so it confounds session diversity with data quantity.

\paragraph{Bipolar re-referencing} To use a bipolar pair as the selection unit, we formed $B_j = F_j - F_{j+8}$ ($j = 1, \ldots, 8$) from the monopolar recordings, re-extracted features with the threshold given by the same rule ($4 \times 10^{-3}$~mV; the bipolar signal is about half the monopolar amplitude), and repeated the analysis with the $k = 8$ model as reference. These are derived pairs, not physical bipolar electrodes.

\subsection{Departures from the analysis plan}
\label{sec:deviations}
The detailed H2--H4 test families were specified after the underlying results were available; the plan's Holm correction for H1 was applied after the primary analysis; $\lambda$ was not selected; training-sample sizes were chosen on two folds; the putEMG replication, a fixed-layout check and a feature-redundancy check were added; and two exploratory analyses listed in the plan were not run. None of these changes altered a prespecified decision rule. The complete list is given in supplementary methods~S1.
